(a) Take $a=1,b=9,n=2,k=1$. Then $n\mid a+b$, but the proposed divisor is $2^5=32$, which does not divide $1-81=-80$. Also the original sign may fail even divisibility by $n$ when $n$ is odd and $k$ is even; for instance $a=1,b=2,n=3,k=2$ gives $65$.

(b) Write $a=-b+nt$ for an integer $t$. Modulo $n^2$, the binomial expansion has only its constant and linear terms:
\[
a^{nk}\equiv(-b)^{nk}+nk(-b)^{nk-1}nt
\equiv(-1)^{nk}b^{nk}\pmod{n^2}.
\]
Every term of order at least two already contains $n^2$, and the linear term does as well. This proves the corrected assertion, including the trivial case $n=1$.

(c) Use $a=6,b=8,n=7,k=39$. Then $7\mid14$ and $nk=273$ is odd, so the corrected statement gives
\[
\boxed{6^{273}+8^{273}\equiv0\pmod{49}}.
\]
